\chapter*{Введение}
\addcontentsline{toc}{chapter}{Введение}
\section*{Актуальность}
Амплитуда регистрируемого вызванного составного потенциала действия (ECAP) зависит от параметров стимуляции и условий измерения. Моделирование позволяет исследовать вклад возбуждения волокон, геометрии и регистрации в наблюдаемый ответ~\cite{ecapmodel},~\cite{posture}.

\section*{Цель исследования}
Разработать метод биофизического моделирования ECAP с предобучением динамического представления на коннектомной модели Drosophila и установить наличие, отсутствие либо неопределённость его дополнительного вклада в точность моделирования кривой роста ECAP на независимых данных человека.
\section*{Объект и предмет}
Объект --- процесс формирования и регистрации ECAP в измерительном канале систем стимуляции спинного мозга (SCS). Предмет --- биофизическое моделирование ECAP, предобучение и адаптация временного модуля, оценка неопределённости и условий применимости.
\section*{Задачи исследования}
\begin{enumerate}
\setlength{\itemsep}{3pt}

\item Проанализировать методы коннектомного моделирования, предобучения и биофизического моделирования ECAP. Определить требования к происхождению, совместимости и условиям использования исходных данных; выбрать основной и независимый наборы ECAP после подтверждения доступа и пригодности для исследования.
\item Разработать коннектомно-ограниченную модель нейродинамики Drosophila и алгоритм генерации синтетических временных рядов. Проверить модель по наблюдениям, совместимым с выбранной версией коннектома и стадией развития; оценить влияние структуры связей и параметров динамики на воспроизводимые характеристики активности.
\item Определить состав переносимого временного представления и разработать алгоритм предобучения и адаптации начальных параметров целевого модуля. Проверить допустимость сопоставления динамических характеристик и их значение для моделирования формы ECAP и кривой роста без предположения анатомической эквивалентности видов.
\item Разработать прямой биофизический оператор формирования и регистрации ECAP с учётом рекрутирования волокон, распространения возбуждения, объёмной проводимости, геометрии электродов, артефакта и шума. Провести его отдельную проверку на независимых человеческих записях при подтверждённой пригодности и доступности данных.
\item Разработать и заранее зафиксировать протокол сравнительной оценки метода переноса, включая контрольные модели, основную метрику, допуски, разбиения и статистический анализ. Провести сравнение на участниках, не использованных при обучении и выборе модели; оценить вклад структуры, динамики и предобучения, а также необходимость временного модуля.
\item Установить область применимости метода, чувствительность к параметрам измерительного канала и условия получения достоверных оценок неопределённости. Разработать программную реализацию и методику воспроизведения расчётов; определить границы технических выводов для систем стимуляции спинного мозга.
\end{enumerate}
\section*{Проверяемая гипотеза}
Предобучение временного модуля на коннектомно-ограниченной динамике Drosophila может снизить ошибку моделирования кривой роста ECAP относительно обучения без переноса и контрольных вариантов. Предполагается проверить практически значимое преимущество, его подтверждённое отсутствие, ухудшение либо неопределённость при сопоставимых данных, архитектуре и вычислительных затратах.
\section*{Предполагаемый вклад}
Планируется разработать способ сопряжения переносимого временного модуля с отдельным биофизическим оператором и установить условия его применимости.
